\section{Задание 2. Ортогональная проекция}

1. Найти ортогональную проекцию и ортогональную составляющую вектора $x$ при проекции на подпространство $L=\langle a_1,a_2,a_3\rangle$:

\hspace*{\parindent}1.1. разложив вектор $x$ по базису пространства, состоящему из 

векторов базиса $L$ и базиса $L^\perp$;

\hspace*{\parindent}1.2. используя матрицу Грама базиса $L$;

\hspace*{\parindent}1.3. используя формулу ортогональной проекции.

2. Найти расстояние и угол между вектором $x$ и подпространством $L$.
\vspace{0.1cm}
\[
17.\quad
a_1 =
\begin{pmatrix}
4\\
1\\
3\\
0
\end{pmatrix},\quad
a_2 =
\begin{pmatrix}
-1\\
1\\
-1\\
-1
\end{pmatrix},\quad
a_3 =
\begin{pmatrix}
6\\
-1\\
5\\
2
\end{pmatrix},\quad
x =
\begin{pmatrix}
-1\\
3\\
11\\
3
\end{pmatrix}.
\]

\textbf{Решение.}

\textbf{1. Найдём ортогональную проекцию и ортогональную составляющую вектора $x$ при проекции на подпространство}
\[ 
L=\langle a_1,a_2,a_3\rangle,
\]

\textbf{1.1. Разложение по базису из векторов $L$ и $L^\perp$.}

Найдём базис $L$:
\[
\begin{pmatrix}
4 & -1 & 6\\
1 & 1 & -1\\
3 & -1 & 5\\
0 & -1 & 2
\end{pmatrix}
\sim
\begin{pmatrix}
1 & 1 & -1\\
0 & -5 & 10\\
0 & -4 & 8\\
0 & -1 & 2
\end{pmatrix}
\sim
\begin{pmatrix}
1 & 1 & -1\\
0 & 1 & -2\\
0 & 0 & 0\\
0 & 0 & 0
\end{pmatrix}.
\]

\[
r=2,
\]
следовательно, базисом $L$ являются векторы
\[
a_1=
\begin{pmatrix}
4\\
1\\
3\\
0
\end{pmatrix},
\qquad
a_2=
\begin{pmatrix}
-1\\
1\\
-1\\
-1
\end{pmatrix}.
\]

Найдём базис $L^\perp$.

\[
y\in L^\perp
\iff
(\forall a\in L)\ (a,y)=0
\iff
\begin{cases}
(a_1,y)=0,\\
(a_2,y)=0.
\end{cases}
\]

Пусть
\[
y=
\begin{pmatrix}
y_1\\
y_2\\
y_3\\
y_4
\end{pmatrix}.
\]
Тогда
\[
\begin{cases}
4y_1+y_2+3y_3=0,\\
-y_1+y_2-y_3-y_4=0.
\end{cases}
\]

Выразим переменные $y_1$ и $y_2$ через $y_3$ и $y_4$.

Из второго уравнения получаем:
\[
-y_1+y_2-y_3-y_4=0
\;\Rightarrow\;
y_2=y_1+y_3+y_4.
\]

Подставим это выражение в первое уравнение:
\[
4y_1+(y_1+y_3+y_4)+3y_3=0.
\]

\[
5y_1+4y_3+y_4=0
\;\Rightarrow\;
y_1=-\frac{4}{5}y_3-\frac{1}{5}y_4.
\]

Теперь найдём $y_2$:
\[
y_2=y_1+y_3+y_4=
\left(-\frac{4}{5}y_3-\frac{1}{5}y_4\right)+y_3+y_4.
\]

\[
y_2=\frac{1}{5}y_3+\frac{4}{5}y_4.
\]

Итак,
\[
\begin{cases}
y_1=-\dfrac45y_3-\dfrac15y_4,\\[2mm]
y_2=\dfrac15y_3+\dfrac45y_4,\\
y_3 \text{ --- свободная},\\
y_4 \text{ --- свободная}.
\end{cases}
\]

Положим
\[
y_3=\alpha,\qquad y_4=\beta.
\]

Тогда
\[
y_1=-\frac45\alpha-\frac15\beta,\qquad
y_2=\frac15\alpha+\frac45\beta.
\]

Следовательно,
\[
y=
\begin{pmatrix}
y_1\\
y_2\\
y_3\\
y_4
\end{pmatrix}
=
\begin{pmatrix}
-\frac45\alpha-\frac15\beta\\[1mm]
\frac15\alpha+\frac45\beta\\[1mm]
\alpha\\
\beta
\end{pmatrix}.
\]

Вынесем отдельно коэффициенты при $\alpha$ и $\beta$:
\[
y=
\alpha
\begin{pmatrix}
-\frac45\\[1mm]
\frac15\\[1mm]
1\\
0
\end{pmatrix}
+
\beta
\begin{pmatrix}
-\frac15\\[1mm]
\frac45\\[1mm]
0\\
1
\end{pmatrix}
=
\alpha
\begin{pmatrix}
-4\\[1mm]
1\\[1mm]
5\\
0
\end{pmatrix}
+
\beta
\begin{pmatrix}
-1\\[1mm]
4\\[1mm]
0\\
5
\end{pmatrix}.
\]

Следовательно, базис $L^\perp$ имеет вид
\[
\begin{pmatrix}
-4\\
1\\
5\\
0
\end{pmatrix},
\qquad
\begin{pmatrix}
-1\\
4\\
0\\
5
\end{pmatrix}.
\]

Так как $V=L+L^\perp$, составим базис пространства $V$ из векторов базиса $L$ и базиса $L^\perp$.
Найдём координаты вектора $x$ в этом базисе:
\[
\alpha_1
\begin{pmatrix}
4\\
1\\
3\\
0
\end{pmatrix}
+
\alpha_2
\begin{pmatrix}
-1\\
1\\
-1\\
-1
\end{pmatrix}
+
\alpha_3
\begin{pmatrix}
-4\\
1\\
5\\
0
\end{pmatrix}
+
\alpha_4
\begin{pmatrix}
-1\\
4\\
0\\
5
\end{pmatrix}
=
\begin{pmatrix}
-1\\
3\\
11\\
3
\end{pmatrix}.
\]

\[
\left(
\begin{array}{cccc|c}
4 & -1 & -4 & -1 & -1\\
1 & 1 & 1 & 4 & 3\\
3 & -1 & 5 & 0 & 11\\
0 & -1 & 0 & 5 & 3
\end{array}
\right)
\sim
\left(
\begin{array}{cccc|c}
1 & 1 & 1 & 4 & 3\\
4 & -1 & -4 & -1 & -1\\
3 & -1 & 5 & 0 & 11\\
0 & -1 & 0 & 5 & 3
\end{array}
\right)
\]

\[
\sim
\left(
\begin{array}{cccc|c}
1 & 1 & 1 & 4 & 3\\
0 & -5 & -8 & -17 & -13\\
0 & -4 & 2 & -12 & 2\\
0 & -1 & 0 & 5 & 3
\end{array}
\right)
\sim
\left(
\begin{array}{cccc|c}
1 & 1 & 1 & 4 & 3\\
0 & -1 & 0 & 5 & 3\\
0 & -4 & 2 & -12 & 2\\
0 & -5 & -8 & -17 & -13
\end{array}
\right)
\]

\[
\sim
\left(
\begin{array}{cccc|c}
1 & 1 & 1 & 4 & 3\\
0 & 1 & 0 & -5 & -3\\
0 & -4 & 2 & -12 & 2\\
0 & -5 & -8 & -17 & -13
\end{array}
\right)
\sim
\left(
\begin{array}{cccc|c}
1 & 0 & 1 & 9 & 6\\
0 & 1 & 0 & -5 & -3\\
0 & 0 & 2 & -32 & -10\\
0 & 0 & -8 & -42 & -28
\end{array}
\right)
\]

\[
\sim
\left(
\begin{array}{cccc|c}
1 & 0 & 1 & 9 & 6\\
0 & 1 & 0 & -5 & -3\\
0 & 0 & 1 & -16 & -5\\
0 & 0 & -8 & -42 & -28
\end{array}
\right)
\sim
\left(
\begin{array}{cccc|c}
1 & 0 & 0 & 25 & 11\\
0 & 1 & 0 & -5 & -3\\
0 & 0 & 1 & -16 & -5\\
0 & 0 & 0 & -170 & -68
\end{array}
\right)
\]

\[
\sim
\left(
\begin{array}{cccc|c}
1 & 0 & 0 & 25 & 11\\
0 & 1 & 0 & -5 & -3\\
0 & 0 & 1 & -16 & -5\\
0 & 0 & 0 & 1 & \frac25
\end{array}
\right)
\sim
\left(
\begin{array}{cccc|c}
1 & 0 & 0 & 0 & 1\\
0 & 1 & 0 & 0 & -1\\
0 & 0 & 1 & 0 & \frac75\\
0 & 0 & 0 & 1 & \frac25
\end{array}
\right).
\]

Отсюда
\[
\alpha_1=1,\qquad
\alpha_2=-1,\qquad
\alpha_3=\frac75,\qquad
\alpha_4=\frac25.
\]

Следовательно,
\[
x=
\left(
\begin{pmatrix}
4\\
1\\
3\\
0
\end{pmatrix}
-
\begin{pmatrix}
-1\\
1\\
-1\\
-1
\end{pmatrix}
\right)
+
\left(
\frac75
\begin{pmatrix}
-4\\
1\\
5\\
0
\end{pmatrix}
+
\frac25
\begin{pmatrix}
-1\\
4\\
0\\
5
\end{pmatrix}
\right).
\]

Значит,
\[
x_L=
\begin{pmatrix}
5\\
0\\
4\\
1
\end{pmatrix},
\qquad
x^\perp=
\begin{pmatrix}
-6\\
3\\
7\\
2
\end{pmatrix}.
\]

\bigskip

\textbf{1.2. Использование матрицы Грама.}

$L$ --- подпространство $V$. Векторы $a_1,a_2$ образуют базис $L$.

\[
(\forall x\in V)\quad x=x_L+x^\perp,\qquad x_L\in L,\quad x^\perp\in L^\perp.
\]

Пусть
\[
x_L=\alpha_1a_1+\alpha_2a_2.
\]

Умножим скалярно на $a_1$ и $a_2$. Получим:
\[
(x_L,a_1)=\alpha_1(a_1,a_1)+\alpha_2(a_2,a_1),
\]
\[
(x_L,a_2)=\alpha_1(a_1,a_2)+\alpha_2(a_2,a_2).
\]

Так как $x=x_L+x^\perp$ и $(x^\perp,a_1)=0$, $(x^\perp,a_2)=0$, то
\[
(x_L,a_1)=(x,a_1),\qquad (x_L,a_2)=(x,a_2).
\]

Получаем:
\[
G(a_1,a_2)
\begin{pmatrix}
\alpha_1\\
\alpha_2
\end{pmatrix}
=
\begin{pmatrix}
(x,a_1)\\
(x,a_2)
\end{pmatrix}.
\]

Вычислим:
\[
(a_1,a_1)=4^2+1^2+3^2=26,
\]
\[
(a_1,a_2)=4\cdot(-1)+1\cdot1+3\cdot(-1)+0\cdot(-1)=-6,
\]
\[
(a_2,a_2)=(-1)^2+1^2+(-1)^2+(-1)^2=4.
\]

Следовательно,
\[
G(a_1,a_2)=
\begin{pmatrix}
26 & -6\\
-6 & 4
\end{pmatrix}.
\]

Далее
\[
(x,a_1)=(-1)\cdot4+3\cdot1+11\cdot3+3\cdot0=32,
\]
\[
(x,a_2)=(-1)\cdot(-1)+3\cdot1+11\cdot(-1)+3\cdot(-1)=-10.
\]

Получаем систему
\[
\begin{pmatrix}
26 & -6\\
-6 & 4
\end{pmatrix}
\begin{pmatrix}
\alpha_1\\
\alpha_2
\end{pmatrix}
=
\begin{pmatrix}
32\\
-10
\end{pmatrix},
\]
то есть
\[
\begin{cases}
26\alpha_1-6\alpha_2=32,\\
-6\alpha_1+4\alpha_2=-10.
\end{cases}
\]

Разделим второе уравнение на $2$:
\[
\begin{cases}
26\alpha_1-6\alpha_2=32,\\
-3\alpha_1+2\alpha_2=-5.
\end{cases}
\]

Умножим второе уравнение на $3$:
\[
\begin{cases}
26\alpha_1-6\alpha_2=32,\\
-9\alpha_1+6\alpha_2=-15.
\end{cases}
\]

Сложим уравнения:
\[
17\alpha_1=17,
\]
откуда
\[
\alpha_1=1.
\]

Подставим в уравнение
\[
-3\alpha_1+2\alpha_2=-5:
\]
\[
-3+2\alpha_2=-5,
\]
\[
2\alpha_2=-2,
\]
\[
\alpha_2=-1.
\]

Следовательно,
\[
x_L=\alpha_1a_1+\alpha_2a_2=
\begin{pmatrix}
4\\
1\\
3\\
0
\end{pmatrix}
-
\begin{pmatrix}
-1\\
1\\
-1\\
-1
\end{pmatrix}
=
\begin{pmatrix}
5\\
0\\
4\\
1
\end{pmatrix}.
\]

Тогда
\[
x^\perp=x-x_L=
\begin{pmatrix}
-1\\
3\\
11\\
3
\end{pmatrix}
-
\begin{pmatrix}
5\\
0\\
4\\
1
\end{pmatrix}
=
\begin{pmatrix}
-6\\
3\\
7\\
2
\end{pmatrix}.
\]

\bigskip

\textbf{1.3. Использование формулы ортогональной проекции.}

Базис $L$:
\[
a_1=
\begin{pmatrix}
4\\
1\\
3\\
0
\end{pmatrix},
\qquad
a_2=
\begin{pmatrix}
-1\\
1\\
-1\\
-1
\end{pmatrix}.
\]

Найдём ортогональный базис подпространства $L$ при помощи процесса ортогонализации Грама--Шмидта.

\[
b_1=a_1=
\begin{pmatrix}
4\\
1\\
3\\
0
\end{pmatrix}.
\]

\[
b_2=a_2-\frac{(b_1,a_2)}{(b_1,b_1)}\,b_1.
\]

Вычислим скалярные произведения:
\[
(b_1,a_2)=4\cdot(-1)+1\cdot1+3\cdot(-1)+0\cdot(-1)=-6,
\]
\[
(b_1,b_1)=4^2+1^2+3^2+0^2=16+1+9=26.
\]

Тогда
\[
b_2=
\begin{pmatrix}
-1\\
1\\
-1\\
-1
\end{pmatrix}
+\frac{6}{26}
\begin{pmatrix}
4\\
1\\
3\\
0
\end{pmatrix}
=
\begin{pmatrix}
-1\\
1\\
-1\\
-1
\end{pmatrix}
+\frac{3}{13}
\begin{pmatrix}
4\\
1\\
3\\
0
\end{pmatrix}
=
\begin{pmatrix}
-\frac{1}{13}\\
\frac{16}{13}\\
-\frac{4}{13}\\
-1
\end{pmatrix}
=
\frac{1}{13}
\begin{pmatrix}
-1\\
16\\
-4\\
-13
\end{pmatrix}.
\]

Поэтому можно взять
\[
b_2=
\begin{pmatrix}
-1\\
16\\
-4\\
-13
\end{pmatrix}.
\]

Так как $b_1$ и $b_2$ образуют ортогональный базис подпространства $L$, ортогональная проекция вектора $x$ на это подпространство находится по формуле
\[
x_L=
\frac{(b_1,x)}{(b_1,b_1)}\,b_1+
\frac{(b_2,x)}{(b_2,b_2)}\,b_2.
\]

Имеем
\[
(b_1,b_1)=26,\qquad (b_1,x)=-4+3+33+0=32,
\]
\[
(b_2,b_2)=1+256+16+169=442,\qquad (b_2,x)=1+48-44-39=-34.
\]

Следовательно,
\[
x_L=
\frac{32}{26}
\begin{pmatrix}
4\\
1\\
3\\
0
\end{pmatrix}
+
\frac{-34}{442}
\begin{pmatrix}
-1\\
16\\
-4\\
-13
\end{pmatrix}
=
\frac{16}{13}
\begin{pmatrix}
4\\
1\\
3\\
0
\end{pmatrix}
-
\frac{1}{13}
\begin{pmatrix}
-1\\
16\\
-4\\
-13
\end{pmatrix}
=\]
\[=
\begin{pmatrix}
\frac{64}{13}\\
\frac{16}{13}\\
\frac{48}{13}\\
0
\end{pmatrix}
-
\begin{pmatrix}
-\frac{1}{13}\\
\frac{16}{13}\\
-\frac{4}{13}\\
-1
\end{pmatrix}
=
\begin{pmatrix}
\frac{65}{13}\\
0\\
\frac{52}{13}\\
1
\end{pmatrix}
=
\begin{pmatrix}
5\\
0\\
4\\
1
\end{pmatrix}.
\]

Тогда
\[
x^\perp=x-x_L=
\begin{pmatrix}
-1\\
3\\
11\\
3
\end{pmatrix}
-
\begin{pmatrix}
5\\
0\\
4\\
1
\end{pmatrix}
=
\begin{pmatrix}
-6\\
3\\
7\\
2
\end{pmatrix}.
\]

\bigskip

\textbf{2. Найдём расстояние и угол между вектором $x$ и подпространством $L$.}

Расстояние между вектором $x$ и подпространством $L$:
\[
\|x^\perp\|=\sqrt{(-6)^2+3^2+7^2+2^2}=\sqrt{36+9+49+4}=\sqrt{98}=7\sqrt2.
\]

Угол между вектором $x$ и подпространством $L$:
\[
\varphi=\arcsin\frac{\|x^\perp\|}{\|x\|}.
\]
\[
\|x\|=\sqrt{(-1)^2+3^2+11^2+3^2}=\sqrt{1+9+121+9}=\sqrt{140}=2\sqrt{35}.
\]

Следовательно,
\[
\varphi=\arcsin\frac{7\sqrt2}{2\sqrt{35}}
=
\arcsin\sqrt{\frac{98}{140}}
=
\arcsin\sqrt{\frac{7}{10}}.
\]
